\documentclass[11pt,letterpaper]{article}

\usepackage[utf8]{inputenc}
\usepackage[spanish]{babel}
\usepackage[margin=2.5cm]{geometry}
\usepackage{amsmath,amssymb,amsthm}
\usepackage{enumitem}
\usepackage{titlesec}
\usepackage{hyperref}
\usepackage{xcolor}
\usepackage{listings}

\definecolor{unicaucaverde}{RGB}{0,90,60}
\hypersetup{colorlinks=true,linkcolor=unicaucaverde,urlcolor=unicaucaverde}

\titleformat{\section}{\normalfont\Large\bfseries\color{unicaucaverde}}{\thesection}{1em}{}
\titleformat{\subsection}{\normalfont\large\bfseries}{\thesubsection}{1em}{}

\theoremstyle{definition}
\newtheorem{definicion}{Definición}[section]
\newtheorem{ejemplo}[definicion]{Ejemplo}
\newtheorem{propiedad}[definicion]{Propiedad}

\lstdefinestyle{cstyle}{
  language=C,
  basicstyle=\ttfamily\small,
  keywordstyle=\color{unicaucaverde}\bfseries,
  commentstyle=\color{gray}\itshape,
  stringstyle=\color{orange!70!black},
  numbers=left,
  numberstyle=\tiny\color{gray},
  frame=single,
  breaklines=true,
  showstringspaces=false,
  tabsize=2,
  captionpos=b
}
\lstset{style=cstyle}

\title{\textcolor{unicaucaverde}{\Large Notas de clase --- Semana 4}\\[4pt]
  \large Ámbito avanzado y recursión}
\author{Programación --- Universidad del Cauca}
\date{2026}

\begin{document}
\maketitle
\tableofcontents

\section{Variables estáticas locales: tiempo de vida frente a ámbito}

\begin{definicion}[Ámbito y tiempo de vida]
El \emph{ámbito} (\emph{scope}) de una variable es la región del código donde su nombre es visible
y se puede usar. El \emph{tiempo de vida} (\emph{storage duration}) es el período durante el cual
la variable existe realmente en memoria. En una variable local ordinaria ambos coinciden: nace y
muere con cada ejecución de la función. La palabra clave \texttt{static} rompe esa coincidencia.
\end{definicion}

\begin{definicion}[Variable estática local]
Una variable local declarada con el calificador \texttt{static} conserva su \emph{ámbito} normal
(solo es visible dentro de la función donde se declaró), pero su \emph{tiempo de vida} pasa a ser
el de todo el programa: se inicializa una única vez, y conserva su valor de una llamada a la
siguiente en lugar de recrearse cada vez.
\end{definicion}

\begin{ejemplo}[Contador que persiste entre llamadas]
\begin{lstlisting}
#include <stdio.h>

int contadorLlamadas(void) {
    static int veces = 0;   // se inicializa UNA sola vez, en toda la ejecucion
    veces++;
    return veces;
}

int main(void) {
    printf("%d\n", contadorLlamadas());   // imprime 1
    printf("%d\n", contadorLlamadas());   // imprime 2
    printf("%d\n", contadorLlamadas());   // imprime 3
    return 0;
}
\end{lstlisting}
A diferencia de una variable local ordinaria (que reiniciaría \texttt{veces} a 0 en cada llamada),
\texttt{static int veces = 0;} solo ejecuta la inicialización \texttt{= 0} la primera vez que el
programa pasa por esa línea; en las llamadas siguientes conserva el valor que tenía al salir de la
función anterior.
\end{ejemplo}

\begin{propiedad}[\texttt{static} local frente a variable global]
Una variable estática local y una variable global comparten el mismo tiempo de vida (todo el
programa), pero difieren en ámbito: la global es visible desde cualquier función que la vea
declarada, mientras que la estática local solo es visible \emph{dentro} de la función donde se
declaró. Esto la hace preferible cuando el estado que se quiere conservar (un contador, una
bandera de "primera vez") le pertenece conceptualmente a una sola función y no debe poder ser leído
ni modificado por el resto del programa.
\end{propiedad}

\section{Ámbito avanzado: bloques anidados y sombreado}

\begin{definicion}[Bloque]
En C, cualquier par de llaves \texttt{\{ \}} delimita un \emph{bloque}: el cuerpo de una función,
de un \texttt{if}, de un \texttt{for}, o incluso un bloque suelto. Una variable declarada dentro de
un bloque es local a ese bloque: su ámbito termina en la llave de cierre \texttt{\}}
correspondiente.
\end{definicion}

\begin{ejemplo}[Sombreado dentro de bloques anidados]
\begin{lstlisting}
#include <stdio.h>

int main(void) {
    int x = 1;
    printf("Afuera: x = %d\n", x);        // 1
    {
        int x = 2;                        // sombrea (oculta) la x externa
        printf("Bloque interno: x = %d\n", x);  // 2
    }
    printf("Afuera de nuevo: x = %d\n", x);     // 1 (la x externa no cambio)
    return 0;
}
\end{lstlisting}
Dentro del bloque interno, el nombre \texttt{x} se refiere a la nueva variable declarada ahí (que
\emph{sombrea} a la externa); al salir del bloque, esa \texttt{x} interna desaparece y el nombre
vuelve a referirse a la \texttt{x} original, que nunca fue tocada.
\end{ejemplo}

\begin{propiedad}[Regla práctica]
Combinar \texttt{static} con bloques anidados o con nombres repetidos entre funciones distintas es
una fuente común de errores difíciles de rastrear: antes de declarar una variable estática o de
reutilizar un nombre en un bloque anidado, conviene preguntarse explícitamente qué ámbito y qué
tiempo de vida necesita realmente esa variable.
\end{propiedad}

\section{Recursión: caso base y caso recursivo}

\begin{definicion}[Función recursiva]
Una función es \emph{recursiva} cuando, en su propio cuerpo, se llama a sí misma (directamente, o
indirectamente a través de otra función). Toda función recursiva correcta necesita dos partes: un
\emph{caso base} que se resuelve directamente sin volver a llamarse, y uno o más \emph{casos
recursivos} que reducen el problema a una versión más pequeña del mismo problema y se apoyan en esa
llamada más pequeña para construir el resultado.
\end{definicion}

\begin{ejemplo}[Factorial recursivo]
\begin{lstlisting}
#include <stdio.h>

int factorial(int n) {
    if (n <= 1) {           // caso base
        return 1;
    }
    return n * factorial(n - 1);   // caso recursivo
}

int main(void) {
    printf("5! = %d\n", factorial(5));   // 120
    return 0;
}
\end{lstlisting}
\texttt{factorial(5)} llama a \texttt{factorial(4)}, que llama a \texttt{factorial(3)}, y así
sucesivamente hasta \texttt{factorial(1)}, que es el caso base y devuelve 1 sin generar más
llamadas. A partir de ahí, cada llamada pendiente multiplica su \texttt{n} por el resultado que
recibió: $1 \to 1 \to 2 \to 6 \to 24 \to 120$.
\end{ejemplo}

\begin{propiedad}[Por qué el caso base es obligatorio]
Sin un caso base que detenga las llamadas, una función recursiva se llama a sí misma
indefinidamente. Cada llamada consume memoria adicional (ver \S 4), así que una recursión sin caso
base --- o con un caso base que nunca se alcanza para ciertas entradas --- termina agotando esa
memoria (\emph{stack overflow}) y el programa se detiene con un error en tiempo de ejecución, no en
tiempo de compilación.
\end{propiedad}

\section{La pila de llamadas; recursión frente a iteración}

\begin{definicion}[Pila de llamadas]
La \emph{pila de llamadas} (\emph{call stack}) es la estructura que usa el programa en ejecución
para llevar el control de las funciones activas: cada vez que se llama a una función se apila un
nuevo \emph{marco} con sus parámetros y variables locales, y ese marco se retira de la pila
(\emph{desapila}) cuando la función termina y devuelve su resultado a quien la llamó. En una
función recursiva, cada llamada activa agrega un marco nuevo, distinto de los anteriores: los
parámetros \texttt{n} de \texttt{factorial(5)}, \texttt{factorial(4)}, \ldots, \texttt{factorial(1)}
coexisten en la pila, cada uno en su propio marco, mientras las llamadas están pendientes de
terminar.
\end{definicion}

\begin{ejemplo}[Fibonacci: recursivo frente a iterativo]
\begin{lstlisting}
// Version recursiva
int fibonacciRec(int n) {
    if (n == 0) return 0;         // caso base 1
    if (n == 1) return 1;         // caso base 2
    return fibonacciRec(n - 1) + fibonacciRec(n - 2);   // caso recursivo
}

// Version iterativa equivalente
int fibonacciIter(int n) {
    if (n == 0) return 0;
    int anterior = 0, actual = 1;
    for (int i = 2; i <= n; i++) {
        int siguiente = anterior + actual;
        anterior = actual;
        actual = siguiente;
    }
    return actual;
}
\end{lstlisting}
Ambas funciones calculan el mismo resultado, pero \texttt{fibonacciRec} genera dos llamadas nuevas
por cada llamada que no es caso base (crece exponencialmente en número de llamadas), mientras que
\texttt{fibonacciIter} usa un solo ciclo y un tamaño de memoria constante. La versión recursiva
suele ser más clara de leer porque refleja directamente la definición matemática; la versión
iterativa suele ser más eficiente en tiempo y memoria. Elegir entre una y otra es una decisión de
diseño, no una regla fija.
\end{ejemplo}

\begin{propiedad}[Cuándo preferir cada una]
Se prefiere recursión cuando el problema tiene una estructura naturalmente recursiva y la claridad
del código importa más que el costo adicional de las llamadas (por ejemplo, recorrer una
estructura con subestructuras del mismo tipo). Se prefiere iteración cuando el problema es
esencialmente un conteo o acumulación simple, o cuando el número de llamadas recursivas necesarias
podría ser muy grande y el costo en memoria de la pila de llamadas es una preocupación real.
\end{propiedad}

\section{Ejercicios propuestos}

\begin{enumerate}[leftmargin=1.2cm]
  \item Escribe una función \texttt{int siguienteId(void)} que use una variable estática local para
        devolver 1 en la primera llamada, 2 en la segunda, 3 en la tercera, y así sucesivamente.
  \item Escribe una función recursiva \texttt{int potencia(int base, int exp)} que calcule
        $\text{base}^{\text{exp}}$ (con \texttt{exp} $\geq 0$), identificando claramente su caso
        base.
  \item Escribe una función recursiva \texttt{int sumaHasta(int n)} que calcule
        $1 + 2 + \cdots + n$ para $n \geq 1$.
  \item Escribe una función recursiva \texttt{int fibonacci(int n)} (como en el ejemplo de \S 4) y
        una función iterativa \texttt{int fibonacciIter(int n)} equivalente; verifica que ambas
        producen el mismo resultado para $n = 0, 1, \ldots, 10$.
  \item Dado el código
        \texttt{int f(void) \{ static int c = 0; c++; return c; \}}, llamado tres veces seguidas
        desde \texttt{main} e imprimiendo cada resultado, indica qué se imprime y explica por qué.
  \item Explica con tus palabras por qué una función recursiva sin caso base (o con un caso base
        inalcanzable) eventualmente provoca un error de \emph{stack overflow} en vez de un error de
        compilación.
\end{enumerate}

\section{Soluciones de los ejercicios propuestos}

\begin{description}

\item[\textbf{Ejercicio 1.}]
\begin{lstlisting}
#include <stdio.h>

int siguienteId(void);

int main(void) {
    printf("%d\n", siguienteId());
    printf("%d\n", siguienteId());
    printf("%d\n", siguienteId());
    return 0;
}

int siguienteId(void) {
    static int id = 0;
    id++;
    return id;
}
\end{lstlisting}
Imprime \texttt{1}, \texttt{2}, \texttt{3}.

\item[\textbf{Ejercicio 2.}]
\begin{lstlisting}
#include <stdio.h>

int potencia(int base, int exp);

int main(void) {
    printf("2^5 = %d\n", potencia(2, 5));
    return 0;
}

int potencia(int base, int exp) {
    if (exp == 0) {          // caso base
        return 1;
    }
    return base * potencia(base, exp - 1);   // caso recursivo
}
\end{lstlisting}

\item[\textbf{Ejercicio 3.}]
\begin{lstlisting}
#include <stdio.h>

int sumaHasta(int n);

int main(void) {
    printf("Suma 1..6 = %d\n", sumaHasta(6));
    return 0;
}

int sumaHasta(int n) {
    if (n == 1) {             // caso base
        return 1;
    }
    return n + sumaHasta(n - 1);   // caso recursivo
}
\end{lstlisting}

\item[\textbf{Ejercicio 4.}]
\begin{lstlisting}
#include <stdio.h>

int fibonacci(int n) {
    if (n == 0) return 0;
    if (n == 1) return 1;
    return fibonacci(n - 1) + fibonacci(n - 2);
}

int fibonacciIter(int n) {
    if (n == 0) return 0;
    int anterior = 0, actual = 1;
    for (int i = 2; i <= n; i++) {
        int siguiente = anterior + actual;
        anterior = actual;
        actual = siguiente;
    }
    return actual;
}

int main(void) {
    for (int n = 0; n <= 10; n++) {
        printf("n=%d rec=%d iter=%d\n", n, fibonacci(n), fibonacciIter(n));
    }
    return 0;
}
\end{lstlisting}
Para $n = 0, \ldots, 10$ ambas funciones producen la misma secuencia:
$0, 1, 1, 2, 3, 5, 8, 13, 21, 34, 55$.

\item[\textbf{Ejercicio 5.}]
Se imprime \texttt{1}, \texttt{2}, \texttt{3}. La variable \texttt{static int c = 0;} se inicializa
una sola vez, en la primera llamada a \texttt{f}; en cada llamada posterior conserva el valor que
tenía al salir de la llamada anterior, así que \texttt{c++} lo va incrementando de forma acumulada
(1, luego 2, luego 3) en vez de reiniciarse a 0 en cada llamada como haría una variable local
ordinaria.

\item[\textbf{Ejercicio 6.}]
Cada llamada recursiva pendiente agrega un marco nuevo a la pila de llamadas (con sus propios
parámetros y variables locales), y ese marco solo se retira cuando la llamada termina y regresa un
valor. Si nunca se alcanza un caso base, ninguna llamada termina nunca: los marcos se siguen
apilando sin límite hasta agotar la memoria reservada para la pila. Esto ocurre en tiempo de
ejecución (mientras el programa corre y va haciendo las llamadas), no en tiempo de compilación,
porque el compilador no puede saber en general si el caso base se alcanzará o no para una entrada
dada.
\[\boxed{\text{Sin caso base alcanzable, la pila de llamadas crece sin límite: \emph{stack overflow} en tiempo de ejecución.}}\]

\end{description}

\end{document}
